\documentclass[../../main.tex]{subfiles}
\begin{document}

\begin{center}
  {\fontsize{48pt}{58pt}\selectfont \textbf{1991}\normalsize\textbf{年度}}
\end{center}

\begin{table}[htb]
  \centering
  \caption{出題テーマと難易度}
  \begin{tabular}{|c|p{9cm}|c|}
    \hline
    \textbf{問題番号} & \textbf{テーマ} & \textbf{難易度} \\
    \hline
    第1問 & 正四面体を倒す操作と確率の漸化式 & \\
    \hline
    第2問 & 動く点光源がつくる影の面積 & \\
    \hline
    第3問 & 三次方程式の実根とグラフ & \\
    \hline
    第4問 & 三角関数の多項式表示と導関数 & \\
    \hline
    第5問 & 格子点と直線の距離 & \\
    \hline
    第6問 & 面積条件からの関数の決定 & \\
    \hline
  \end{tabular}
\end{table}

\begin{table}[htb]
  \centering
  \caption{解答}
  \renewcommand{\arraystretch}{1.4}
  \begin{tabular}{|c|c|p{9cm}|}
    \hline
    \textbf{問題番号} & \textbf{小問} & \textbf{答え} \\
    \hline
    第1問 & -- & $p_n=\dfrac{1}{4}\left\{1+3\left(-\dfrac{1}{3}\right)^n\right\}$ \\
    \hline
    第2問 & -- & 影の領域は解答の図を参照．面積は$\left\{\pi c^2+4(c+1)(3c+1)\right\}ab$ \\
    \hline
    第3問 & (1) & 最小値は$-3$ \\
    \cline{2-3}
    & (2) & グラフは解答の図を参照 \\
    \hline
    第4問 & (1) & 所定の多項式$p_n(x)$，$q_n(x)$が存在する（証明問題） \\
    \cline{2-3}
    & (2) & $p_n'(x)=nq_{n-1}(x)$，$q_n'(x)=-np_{n-1}(x)$（証明問題） \\
    \hline
    第5問 & -- & $r$の最小値は$\dfrac{1}{2\sqrt{29}}$ \\
    \hline
    第6問 & (1) & $S(x)=x^{-4}$，$f(x)-2f(2x)=4x^{-5}$ \\
    \cline{2-3}
    & (2) & $\displaystyle\int_x^{2x}a(t)\,dt=0$ \\
    \cline{2-3}
    & (3) & $f(x)=\dfrac{64}{15}x^{-5}$ \\
    \hline
  \end{tabular}
\end{table}

\paragraph{試験形態}
6問．試験時間は要確認．

\end{document}
